% Created 2026-05-14 Thu 21:12
% Intended LaTeX compiler: pdflatex
\documentclass[11pt, letterpaper]{article}

\usepackage[utf8]{inputenc}
\usepackage[T1]{fontenc}
\usepackage{graphicx}
\usepackage{longtable}
\usepackage{wrapfig}
\usepackage{rotating}
\usepackage[normalem]{ulem}
\usepackage{amsmath}
\usepackage{amssymb}
\usepackage{capt-of}
\usepackage{hyperref}
\usepackage[margin=2.5cm]{geometry}      % Márgenes decentes
\usepackage[utf8]{inputenc}
\usepackage{palatino}                   % Tipografía elegante
\usepackage{xcolor}                     % Colores personalizados
\usepackage{listings}
\author{Nieto Gallegos Isaac Julián}
\date{\today}
\title{Instalación de Alpine Linux}
\hypersetup{
 pdfauthor={Nieto Gallegos Isaac Julián},
 pdftitle={Instalación de Alpine Linux},
 pdfkeywords={},
 pdfsubject={},
 pdfcreator={},
 pdflang={English}}
\begin{document}

\maketitle
\tableofcontents

\section{Objetivos}
\label{sec:org47e3e54}

Instalar la distribución de Linux elegida. En este caso, Alpine Linux.
\section{Montaje del sistema raíz}
\label{sec:org4b14959}

\textbf{hasta este punto llevo 13:50 min de la run oficial :p}

El mecanismo de instalación de Alpine Linux está constituido por dos scripts alternativa:

\begin{itemize}
\item setup-alpine, el cual puede ser aplicado sobre un disco duro entero y realizará la instalación automáticamente, según las especificaciones dadas por el usuario a través de variables de entorno y preguntas interactivas del mismo script.

\item setup-disk, el cual es aplicado sobre un arbol de archivos ya montado en \textbf{/mnt}. Sobre este sistema se detectan las particiones utilizadas y se realiza la instalación sobre estas mismas.
\end{itemize}

Para nuestro caso de uso, el segundo script será el usado. Por lo que lo primero que debemos de hacer es crear el arbol raíz en /mnt.

Montamos nuestra partición raíz sobre /mnt con:

\begin{verbatim}
mount -t ext4 <particion> /mnt
\end{verbatim}

Luego, dentro de este, creamos la carpeta /boot con:

\begin{verbatim}
mkdir -p /mnt/boot/efi
mount -t vfat <particion> /mnt/boot/efi
\end{verbatim}

Y finalmente, cargamos la partición swap con:

\begin{verbatim}
swapon <particion>
\end{verbatim}

Una vez hecho esto, tenemos preparado el arbol de archivos sobre el que se instalará el sistema operativo.
\section{Instalación de Alpine Linux}
\label{sec:org4704938}

Instalamos los siguientes paquetes por si acaso:

\begin{verbatim}
apk add grub-efi efibootmgr
\end{verbatim}

Corremos el comando:

\begin{verbatim}
setup-disk -m sys /mnt
\end{verbatim}

Por si acaso, al finalizar corremos los siguiente comandos para registrar la entrada de arranque:

\begin{verbatim}
chroot /mnt /bin/sh
grub-install --target=x86_64-efi --efi-directory=/boot/efi --bootloader-id=alpine
grub-mkconfig -o /boot/grub/grub.cfg
\end{verbatim}
\end{document}
